\documentclass[11pt,a4paper]{article}
\usepackage[T1]{fontenc}
\usepackage{lmodern}
\usepackage[margin=25mm]{geometry}
\usepackage{amsmath,amssymb,amsthm,booktabs,microtype,enumitem}
\usepackage[hidelinks]{hyperref}
\hypersetup{pdftitle={Iterate congruences and a phase correction for A396798},pdfauthor={Dev Goyal},pdfsubject={A396798; formal power series and Lean proofs}}
\setlength{\parskip}{3pt}
\setlength{\emergencystretch}{3em}
\newtheorem{theorem}{Theorem}[section]
\newtheorem{lemma}[theorem]{Lemma}
\newtheorem{corollary}[theorem]{Corollary}
\theoremstyle{remark}\newtheorem{remark}[theorem]{Remark}
\newcommand{\Z}{\mathbb Z}
\newcommand{\N}{\mathbb N}
\newcommand{\F}{\mathbb F}
\newcommand{\coeff}[2]{[x^{#1}]#2}
\title{Iterate congruences and a phase correction for A396798}
\author{Dev Goyal}
\date{24 October 2026}
\begin{document}
\maketitle
\begin{center}\fbox{\begin{minipage}{0.91\linewidth}
\textbf{AI disclosure and review status.}
ChatGPT (OpenAI) generated the mathematical arguments, formal proof code,
and this manuscript from Dev Goyal's prompts. Dev Goyal is the sole human
author and initiated the investigation. This consolidated presentation is
AI-generated; independent human review of this presentation is pending.
Source conjectures and prior results retain the credits given below.
\end{minipage}}\end{center}

\begin{abstract}
This paper gives the defining construction and all completed congruence
results for A396798. It includes the full elementary argument for the
consecutive-iterate family, so no companion paper is needed to follow the
proof. The entry is the parameter $r=4$.
The false starting phase of the original coefficient block is explicitly refuted and corrected.
\end{abstract}
\section{Definition, result, and attribution}
For an integer $r\ge1$, let $F_r$ be the normalized series determined by
\begin{equation}\label{eq:family}
 F_r=x+F_r^{\circ r}F_r^{\circ(r+1)},\qquad F_r=x+O(x^2).
\end{equation}
A superscript $\circ k$ means repeated composition; an ordinary power
means multiplication. Set $a_n=[x^n]F_{4}$.
The sequence A396798 is this case $r=4$, with source
indices $n\ge1$ and an additional zero constant coefficient \cite{oeis}.
All congruences below are coefficientwise.

We prove the formula, for every integer $k$,
\[
F_r^{\circ k}\equiv
\frac{x}{1-kx}+
\begin{cases}
\displaystyle\frac{rx^4}{(1-x)^3},&r\text{ even and }k\text{ odd},\\[5pt]
0,&\text{otherwise}
\end{cases}
\pmod{2r}.
\]
The reduced series has exact order $2r$ under composition. Existence,
uniqueness, and the meaning of negative iterates are proved below.

Hanna recorded the source sequences. Sundaram proved the $r=3$ case,
including integer iterates, in September 2026 \cite{sundaram}. The
odd-parameter argument extends that comparison-series method.
The even-parameter argument needs the displayed correction. The source
conjecture labels in the retained collection do not establish exclusive
priority or worldwide openness.
\section{Formal power series and uniqueness}\label{sec:formal}

\subsection{Definitions}

For a commutative ring $R$ with identity, a formal power series is an
expression $U(x)=\sum_{n\ge0}u_nx^n$ with coefficients in $R$.
The notation $R[[x]]$ denotes the set of these series. Equality means
equality of every coefficient. No numerical value is assigned to $x$,
so analytic convergence is not involved.

Addition is coefficientwise, and multiplication is defined by
\[
\coeff{n}{(UV)}=\sum_{j=0}^n u_jv_{n-j}.
\]
When $V(0)=0$, substitution is defined by
$U(V(x))=\sum_{j\ge0}u_jV(x)^j$. Every coefficient receives only
finitely many contributions. Associativity and the distributive laws
then follow coefficient by coefficient from polynomial identities.

A series is \emph{normalized} if its constant coefficient is zero
and its coefficient of $x$ is one. Thus it has the form $x+O(x^2)$.
The notation $O(x^N)$ denotes a series with zero coefficients in
degrees below $N$. In particular, $U-V=O(x^N)$ says that $U$ and $V$
agree in those degrees. Products satisfy
$O(x^a)O(x^b)=O(x^{a+b})$.

For a normalized $U$, define $U^{\circ0}=x$ and
$U^{\circ(k+1)}=U\circ U^{\circ k}$ for $k\ge0$.
Each iterate is normalized. If $U-V=O(x^N)$, then
\begin{equation}\label{eq:iterate-congruence}
U^{\circ k}-V^{\circ k}=O(x^N)\qquad(k\ge0).
\end{equation}
Indeed, this follows by induction on $k$: split a difference of
compositions into the change in the outer series and the change in
the inner series. Substitution by a normalized series preserves the
degree of vanishing, and factoring $P^j-Q^j$ shows that a change of
order $N$ in the inner argument creates no term of smaller degree.

Reduction of integer coefficients modulo an integer commutes with
these operations, because each coefficient is a finite sum of
products. A series with constant coefficient one has a multiplicative
inverse: if $W(0)=0$, then
\[
(1+W)^{-1}=\sum_{j\ge0}(-W)^j.
\]
This is well defined coefficientwise, and successive terms cancel
upon multiplication. It justifies every fraction used below, whose
denominator has constant coefficient one.

\subsection{The canonical solution}

\begin{lemma}\label{lem:unique}
For every commutative ring $R$ with identity and every integer $r\ge1$,
there is exactly one normalized series $F\in R[[x]]$ satisfying
\[
F=x+F^{\circ r}F^{\circ(r+1)}.
\]
\end{lemma}

\begin{proof}
For normalized $U$, set
$\Phi_r(U)=x+U^{\circ r}U^{\circ(r+1)}$.
This series is normalized. Suppose $U-V=O(x^N)$, with $N\ge2$.
By \eqref{eq:iterate-congruence}, corresponding iterates differ by
$O(x^N)$. Each iterate starts with $x$, so
\begin{align*}
\Phi_r(U)-\Phi_r(V)
={}&(U^{\circ r}-V^{\circ r})U^{\circ(r+1)}\\
&+V^{\circ r}(U^{\circ(r+1)}-V^{\circ(r+1)})
=O(x^{N+1}).
\end{align*}
Thus applying $\Phi_r$ improves agreement by one degree.

For existence, start with $U_0=x$ and define
$U_{j+1}=\Phi_r(U_j)$. We have $U_1-U_0=O(x^2)$; the preceding
calculation gives $U_{j+1}-U_j=O(x^{j+2})$ by induction. Hence every
coefficient eventually stops changing. Define $F$ by those stabilized
coefficients. Each coefficient of $\Phi_r(U)$ depends on only finitely
many coefficients of $U$, so the limit satisfies $F=\Phi_r(F)$.

For uniqueness, two normalized solutions agree below degree two.
Whenever they agree below degree $N$, the displayed calculation and
their defining equations force agreement below degree $N+1$.
Induction gives agreement in every degree.
\end{proof}

The proof uses no division, so it applies to residue rings even when
the modulus is composite. In particular it constructs the integer
series $F_r$ in \eqref{eq:family}. Its image under any ring
homomorphism is the canonical solution over the target ring, by
uniqueness. We can therefore identify its reduction modulo $2r$ by
finding any normalized solution of the reduced equation.

The first coefficients are useful for checking the indexing:
\begin{equation}\label{eq:firstcoeff}
F_r(x)=x+x^2+(2r+1)x^3+O(x^4).
\end{equation}
The product on the right side of \eqref{eq:family} starts with $x^2$,
so the second coefficient is one. It follows by induction that
$F_r^{\circ j}=x+jx^2+O(x^3)$. Multiplying the iterates for $j=r$
and $j=r+1$ gives coefficient $2r+1$ in degree three.

\section{Comparison series and a substitution identity}

Define, for an integer $c$,
\begin{equation}\label{eq:tH}
T_c(x)=\frac{x}{1-cx},\qquad T=T_1,\qquad
H(x)=\frac{x^4}{(1-x)^3}.
\end{equation}
These are integer power series. Direct substitution gives
\begin{equation}\label{eq:t-group}
T_c\circ T_d=T_{c+d}.
\end{equation}
Consequently $T^{\circ k}=T_k$ for every nonnegative integer $k$.
The coefficient of $x^n$ in $T_k$ is $k^{n-1}$ for $n\ge2$,
and its coefficient of $x$ is one.

The formal derivative is defined algebraically by
$U'(x)=\sum_{n\ge1}nu_nx^{n-1}$. It satisfies the product and
substitution rules, which can be proved first for monomials and then
coefficientwise. We need the following rule over a ring in which a
particular element has square zero.

\begin{lemma}\label{lem:taylor}
Let $\varepsilon\in R$ satisfy $\varepsilon^2=0$, and let
$u,v\in xR[[x]]$. For every $S\in R[[x]]$,
\[
S(u+\varepsilon v)=S(u)+\varepsilon vS'(u).
\]
\end{lemma}

\begin{proof}
For every nonnegative integer $j$, the binomial theorem gives
$(u+\varepsilon v)^j=u^j+j\varepsilon vu^{j-1}$, with the second
term interpreted as zero for $j=0$. Terms containing two or more
factors of $\varepsilon$ vanish. Multiply by the coefficient of
$x^j$ in $S$ and sum. All substitutions are well defined because
their inner arguments have zero constant coefficient.
\end{proof}

There is no division by a factorial in this rule. It therefore works
in composite-modulus rings.

\begin{lemma}\label{lem:identities}
The following identities hold in $\Z[[x]]$:
\begin{align}
H(x)T'(T(x))+H(T(x))&=\frac{2x^4}{(1-2x)^3},
\label{eq:identity1}\\
(1-x)H(x)-x^2T(x)-x^3T'(x)&=-\frac{2x^3}{1-x}.
\label{eq:identity2}
\end{align}
\end{lemma}

\begin{proof}
We have $T'(x)=(1-x)^{-2}$, so
\[
H(x)T'(T(x))=\frac{x^4}{(1-x)(1-2x)^2},\qquad
H(T(x))=\frac{x^4}{(1-x)(1-2x)^3}.
\]
Their sum has numerator $x^4((1-2x)+1)=2x^4(1-x)$ over the
common denominator $(1-x)(1-2x)^3$. This proves
\eqref{eq:identity1}. For the second identity, use denominator
$(1-x)^2$ to obtain the numerator
\[
x^4-x^3(1-x)-x^3=-2x^3(1-x).
\]
Cancellation gives \eqref{eq:identity2}.
\end{proof}

Finally, expanding three geometric series shows that
\begin{equation}\label{eq:Hcoeff}
\coeff{n}{H}=
\begin{cases}
0,&n<4,\\
\binom{n-2}{2},&n\ge4.
\end{cases}
\end{equation}
The coefficient of $x^d$ in $(1-x)^{-3}$ counts triples of
nonnegative integers with sum $d$. For each first entry $i$ there
are $d-i+1$ choices for the other two, giving
$(d+1)(d+2)/2=\binom{d+2}{2}$. This proves
\eqref{eq:Hcoeff} without an analytic expansion.

\section{The general congruence}

\begin{theorem}\label{thm:general}
Let $r\ge1$ and $k\ge0$ be integers. Then
\begin{equation}\label{eq:general-congruence}
F_r^{\circ k}(x)\equiv
\begin{cases}
\displaystyle\frac{x}{1-kx},&r\text{ odd or }k\text{ even},\\[7pt]
\displaystyle\frac{x}{1-kx}+\frac{rx^4}{(1-x)^3},
&r\text{ even and }k\text{ odd},
\end{cases}
\pmod{2r}.
\end{equation}
\end{theorem}

\begin{proof}
First suppose $r$ is odd. Since $r+1$ is even,
\[
(1-rx)(1-(r+1)x)
=1-(2r+1)x+r(r+1)x^2\equiv1-x\pmod{2r}.
\]
Using \eqref{eq:t-group}, we obtain
\[
x+T^{\circ r}T^{\circ(r+1)}
\equiv x+\frac{x^2}{1-x}=T\pmod{2r}.
\]
Thus $T$ is the normalized solution of the reduced equation.
Lemma~\ref{lem:unique} gives $F_r\equiv T\pmod{2r}$, and
composition yields $F_r^{\circ k}\equiv T_k$ for all $k\ge0$.

Now suppose $r$ is even. Work in the coefficient ring
$R=\Z/(2r)\Z$, and set $f=T+rH$. The image of $r$ in this ring
has square zero, because $r^2$ is a multiple of $2r$.
Lemma~\ref{lem:taylor} and \eqref{eq:identity1} give
\begin{align*}
f^{\circ2}
&=T(T+rH)+rH(T+rH)\\
&=T^{\circ2}+r\bigl(HT'(T)+H(T)\bigr)\\
&=T^{\circ2}+\frac{2rx^4}{(1-2x)^3}
=T^{\circ2}\quad\text{in }R[[x]].
\end{align*}
Because $r$ is even, repeated composition of this identity gives
\[
f^{\circ r}=T^{\circ r}=\frac{x}{1-rx}=x+rx^2
\quad\text{in }R[[x]].
\]
The last equality again uses $r^2=0$ in $R$. Hence
\[
f^{\circ(r+1)}=f(x+rx^2)=f+rx^2f'.
\]
We can now verify the defining equation. After expanding and discarding
terms with a factor of $r^2$, its residual is
\begin{align*}
f-\bigl(x+f^{\circ r}f^{\circ(r+1)}\bigr)
&=(1-x)f-x-rx^2f-rx^3f'\\
&=r\bigl((1-x)H-x^2T-x^3T'\bigr)\\
&=-\frac{2rx^3}{1-x}=0\quad\text{in }R[[x]],
\end{align*}
where \eqref{eq:identity2} proves the penultimate equality.
By uniqueness, the reduction of $F_r$ is $f$.

For $k=2j$, the equality $f^{\circ2}=T^{\circ2}$ gives
$f^{\circ k}=T^{\circ k}=T_k$. For $k=2j+1$, it gives
\[
f^{\circ k}=f\circ T_{2j}=T_k+rH(T_{2j}).
\]
The integer series $T_{2j}$ is congruent to $x$ modulo two.
Therefore $H(T_{2j})-H(x)$ has even integer coefficients, and
multiplication by $r$ makes it zero modulo $2r$. Thus
$f^{\circ k}=T_k+rH$ in $R[[x]]$, as required.
\end{proof}

\begin{corollary}\label{cor:coeff}
Write $h_n=\coeff{n}{H}$ as in \eqref{eq:Hcoeff}, and let
$\epsilon(r,k)$ equal one if $r$ is even and $k$ is odd, and zero
otherwise. For every $r\ge1$, $k\ge0$, and $n\ge2$,
\[
\coeff{n}{F_r^{\circ k}}
\equiv k^{n-1}+r\epsilon(r,k)h_n\pmod{2r}.
\]
The coefficient of degree zero is zero, and that of degree one is one.
Every positive-index coefficient of $F_r$ is odd.
\end{corollary}

\begin{proof}
The coefficient formula follows from the geometric series for $T_k$
and \eqref{eq:Hcoeff}. At $k=1$, the correction is absent for odd
$r$ and is a multiple of the even integer $r$ for even $r$.
Since $2$ divides $2r$, all positive-index coefficients are odd.
\end{proof}

\section{Integer iterates and exact composition order}\label{sec:integer}

We also prove the integer-iterate extension explicitly, with negative
indices defined through the normalized compositional inverse.

\begin{lemma}\label{lem:inverse}
Every normalized series over a commutative ring has a unique
normalized two-sided inverse under composition.
\end{lemma}

\begin{proof}
For $U=x+\sum_{n\ge2}u_nx^n$, seek
$V=x+\sum_{n\ge2}v_nx^n$ with $U(V)=x$.
The coefficient equation in degree $n$ has the form $v_n$ plus
an expression involving only earlier coefficients of $V$.
Indeed, in $V^j$ for $j\ge2$, a coefficient of degree $n$ uses
only indices at most $n-j+1<n$. Thus every $v_n$ is determined
without division, giving a unique right inverse.
Apply the construction to $V$ to obtain $W$ with $V\circ W=x$.
Associativity gives
$U=U\circ(V\circ W)=(U\circ V)\circ W=W$.
It follows that $V\circ U=x$ too.
\end{proof}

For an integer $k<0$, define $F_r^{\circ k}$ by composing the inverse
$-k$ times. The usual law of integer powers holds in this composition
group. Reducing coefficients commutes with the inverse, because it
preserves both inverse identities and the normalized inverse is unique.

\begin{corollary}\label{cor:order}
The reduction of $F_r$ modulo $2r$ has exact order $2r$ under
composition. Formula~\eqref{eq:general-congruence} remains valid
for every integer $k$, including negative values. In particular,
\[
F_r^{\circ k}\equiv x\pmod{2r}
\quad\Longleftrightarrow\quad 2r\mid k.
\]
\end{corollary}

\begin{proof}
For nonnegative $k$, the coefficient of $x^2$ in the formula of
Theorem~\ref{thm:general} is $k$ modulo $2r$, because $H$ has
no terms below degree four. Equality with $x$ therefore forces
$2r\mid k$. Conversely, if $2r\mid k$, then $k$ is even,
$T_k\equiv x\pmod{2r}$, and the correction term is absent.
This proves the order statement.

For any integer $k$, choose the unique $j$ with
$0\le j<2r$ and $j\equiv k\pmod{2r}$. In the composition group
modulo $2r$, the order statement gives
$F_r^{\circ k}=F_r^{\circ j}$. Also $T_k\equiv T_j\pmod{2r}$,
and $k$ and $j$ have the same parity, because $2r$ is even.
The formula for $j$ is therefore exactly the required formula for
$k$. The same coefficient-of-$x^2$ argument proves the final
equivalence for all integer $k$.
\end{proof}

\section{A396798: seven iterate assertions and a phase correction}

The entry A396798 is the case $r=4$ \cite{oeis}. Let $B=F_4$.
Theorem~\ref{thm:general} gives
\[
B^{\circ k}\equiv T_k+4\,\mathbf1_{\{k\text{ odd}\}}H\pmod8.
\]
For even $k=2,4,6,8$, the geometric series truncates modulo eight
as shown in the table below. For odd $k$, we have $k^2\equiv1\pmod8$.
Using the parity pattern of $h_n$, the coefficients from degree two
repeat $k,1,k+4,5$ modulo eight. These calculations prove all seven
iterate assertions recorded in the entry.

\begin{center}
\begin{tabular}{@{}cl@{}}
\toprule
Iterate & Reduction modulo eight\\
\midrule
$B^{\circ2}$ & $x+2x^2+4x^3$\\
$B^{\circ3}$ & coefficients $3,1,7,5$ repeated from degree two\\
$B^{\circ4}$ & $x+4x^2$\\
$B^{\circ5}$ & coefficients $5,1,1,5$ repeated from degree two\\
$B^{\circ6}$ & $x+6x^2+4x^3$\\
$B^{\circ7}$ & coefficients $7,1,3,5$ repeated from degree two\\
$B^{\circ8}$ & $x$\\
\bottomrule
\end{tabular}
\end{center}

At $k=1$, the correct coefficient block is $1,1,5,5$ beginning
at \emph{degree two}. The entry begins this block at degree one;
that phase is false. Equation~\eqref{eq:firstcoeff} gives
\[
B=x+x^2+9x^3+O(x^4).
\]
Thus its third coefficient has remainder one modulo eight, whereas
the stated phase would require five. The corrected assertion is
\[
\coeff{4j+2}{B}\equiv1,\quad
\coeff{4j+3}{B}\equiv1,\quad
\coeff{4j+4}{B}\equiv5,\quad
\coeff{4j+5}{B}\equiv5\pmod8
\qquad(j\ge0).
\]
It follows for all indices from the general theorem.


\section{Formal proof and scope}
The accompanying entry file is \path{lean/Entries/A396798.lean}.
It states the results with explicit Lean types and supplies proof terms
using the complete shared implementation included in the collection.
The manuscript is self-contained mathematically; the Lean entry file uses
those ordinary library imports rather than duplicating the full foundation.

The entry's named theorem interfaces are listed below. The namespace of
each is \path{FinalCollection.A396798}.
\begin{itemize}[leftmargin=1.5em,itemsep=1pt]
\item \path{source_existence_uniqueness}
\item \path{source_normalized}
\item \path{source_equation}
\item \path{exact_composition_order}
\item \path{a396798_iterate_residue}
\item \path{a396798_corrected_residue}
\item \path{a396798_corrected_period}
\item \path{a396798_counterexample}
\item \path{a396798_literal_period_false}
\item \path{a396798_second_iterate}
\item \path{a396798_fourth_iterate}
\item \path{a396798_sixth_iterate}
\item \path{a396798_eighth_iterate}
\item \path{a396798_third_iterate_period}
\item \path{a396798_fifth_iterate_period}
\item \path{a396798_seventh_iterate_period}
\end{itemize}
The preserved dependencies also contain all supporting constructions,
coefficient identities, map-compatibility results, and inverse proofs.
The collection pins Lean 4.19.0 and Mathlib revision
\path{c44e0c8ee63ca166450922a373c7409c5d26b00b}.
The verification report records fresh compilation and the transitive axiom
audit. Finite numerical checks corroborate the definitions and indexing;
the universally quantified Lean theorems establish the all-index claims.
No assertion of worldwide novelty, OEIS acceptance, publication, or completed
human review follows from these checks. The prose may use elementary
identities that the formal proof establishes by a different algebraic route.

\begin{thebibliography}{9}\raggedright
\bibitem{oeis} P. D. Hanna, \emph{A396798}, The On-Line Encyclopedia of Integer
Sequences. \url{https://oeis.org/A396798}. Source recorded in the retained
proof collection, 24 October 2026.
\bibitem{repo} D. Goyal, \emph{OEIS A396843 proof repository}.
\url{https://github.com/dvgyl/oeis-a396843-proof}.
Reference revision \path{995e55adcbe06b2a913e24f12ff73a7faf362f06}.
\bibitem{sundaram} Shiva Priyan Shanmuga Sundaram,
\emph{A Proof of the Congruence Conjectures for OEIS A396797}, September 2026.
\url{https://oeis.org/A396797/a396797.pdf}.
\end{thebibliography}
\end{document}
